\documentclass[../main.tex]{subfiles}
\begin{document}

\textit{This chapter presents the project management methods the system supports, the algorithms
behind its dependency features, and the technologies used to build it.}

\section{Project management for small software teams}
\subsection{Kanban}
Kanban visualises work as cards moving across columns that represent states, and limits the amount
of work in progress so that the team finishes work before starting more~\cite{anderson2010kanban}.
Its strength is continuous flow: there is no fixed iteration, and a blocked item is visible
immediately. The system uses a fixed four-state workflow:
\[
\text{TODO} \rightarrow \text{IN\_PROGRESS} \rightarrow \text{BLOCKED} \rightarrow \text{DONE} \rightarrow \text{TODO}.
\]
The same cycle is implemented on the client and on the server and must stay identical; a task in
BLOCKED carries a reason, so the board explains \emph{why} work has stopped.

\subsection{Scrum}
Scrum is a framework in which a team delivers work in fixed-length iterations called
\emph{sprints}, usually one to four weeks long~\cite{schwaber2020scrum}. Its artefacts are the
\emph{product backlog} (everything that might be done), the \emph{sprint backlog} (what the team
commits to in this sprint) and the increment. Its events are sprint planning, the daily scrum, the
sprint review and the \emph{sprint retrospective}, in which the team inspects how it worked and
decides what to change. Two measures are standard:
\begin{itemize}
  \item \textbf{Burndown}: the remaining estimated work (story points) per day of a sprint,
        compared with an ideal straight line from the commitment down to zero.
  \item \textbf{Velocity}: the story points completed per sprint, used to size the next
        commitment realistically.
\end{itemize}
The system supports Scrum on top of the Kanban board: a task belongs either to the product backlog or to
one sprint; story points are the task's complexity (S=1, M=2, L=3, XL=5).

\section{Dependency graphs and the critical path}
\subsection{Topological ordering}
If task $A$ depends on task $B$, $B$ must be finished first. The tasks and dependencies form a
directed graph $G=(V,E)$ with an edge $B \rightarrow A$. When $G$ has no cycle (a \emph{directed
acyclic graph}, DAG), a \emph{topological order}---a sequence in which every task appears after all
its prerequisites---exists and can be found with Kahn's algorithm~\cite{kahn1962topological}:
repeatedly take a task with no unfinished prerequisite, output it, and remove its outgoing edges.
The system breaks ties between ready tasks by priority, then due date, then creation time, so the
suggested order puts the most urgent unblocked work first. A cycle (which users can create by
mistake) does not make the algorithm fail: tasks left in a cycle are appended at the end.

\subsection{Critical path}
The Critical Path Method~\cite{kelley1959cpm} finds the longest chain of dependent activities; any
delay on it delays the whole project. The system computes the heaviest chain of \emph{unfinished} tasks
with memoised depth-first search, where the weight of a task $t$ is
\[
w(t) = p(t)\cdot c(t), \qquad
p \in \{1,2,5,10\}\ \text{(LOW..CRITICAL)},\quad
c \in \{1,3,5,8\}\ \text{(S..XL)},
\]
and the weight of a chain is the sum of its tasks. Tasks on this path are flagged on the board and in
the DAG view. A subtle bug found during the project concerned memoisation across a cycle: a value
computed for a walk truncated by a cycle depends on the path taken, so caching it made the result
depend on the order of the input array. Such results are now never cached.

\section{Local-first software}
Local-first software keeps the primary copy of the user's data on the user's device and treats the
server as a synchronisation peer~\cite{kleppmann2019localfirst}. The benefits are instant response
and offline operation; the main difficulty is merging concurrent edits. The system adopts a deliberately
conservative policy: the browser writes every change to IndexedDB first and renders it immediately,
then synchronises with the server; while offline, a \emph{personal} project (one member) stays
editable, whereas a \emph{shared} project becomes read-only, because without a merge algorithm two
people editing offline would silently overwrite each other.

\section{Large language models and structured output}
Large language models based on the Transformer architecture~\cite{vaswani2017attention} can
summarise, extract and plan from natural language. Their output is free text, however, and an
application cannot act on a paragraph. The system therefore asks the model for JSON with a fixed schema
(for example \code{\{went\_well, to\_improve, action\_items\}}), strips Markdown fences, validates and
normalises the result (unknown enum values become a default; dependencies naming a task that does not
exist are dropped), and falls back to a deterministic answer when the output is unusable. The model
used is \code{openai/gpt-oss-120b}~\cite{openai2025gptoss} served by Groq through an
OpenAI-compatible API~\cite{groqdocs}.

\section{Technologies}
\begin{table}[H]
\centering
\caption{Technology stack}\label{tab:stack}
\renewcommand{\arraystretch}{1.25}
\begin{tabular}{|p{3.2cm}|p{4.2cm}|p{6.3cm}|}
\hline
\textbf{Layer} & \textbf{Technology} & \textbf{Purpose} \\ \hline
Frontend & Vue 3, TypeScript, Pinia, Vite, Tailwind CSS & Single-page application; Composition API components and a central store~\cite{vuejs} \\ \hline
Local storage & IndexedDB (idb-keyval) & Local-first copy of each project's tasks \\ \hline
Backend & Python, FastAPI, Pydantic v2 & REST API with typed request/response schemas~\cite{fastapi} \\ \hline
Data & SQLAlchemy 2, Alembic, MySQL 8.4 & ORM, versioned migrations, production database (SQLite for development) \\ \hline
Security & JWT (HS256), bcrypt & Stateless sessions~\cite{rfc7519jwt}, password hashing \\ \hline
AI & Groq API, Ollama (optional) & Cloud LLM, optional local LLM \\ \hline
Testing & pytest, Vitest, Vue Test Utils, Playwright & Unit, component and browser automation \\ \hline
Delivery & Docker Compose, nginx, GitHub Actions, Cloudflare Tunnel & Containers, reverse proxy, CI, public access without opening ports~\cite{cloudflaretunnel} \\ \hline
\end{tabular}
\end{table}

\section{Documentation as a control system}
Documentation that describes a system becomes wrong as soon as the code changes without it; this
project inherited a specification describing behaviour the code did not have. The team therefore
kept eight documents next to the code (D1--D8, listed in the Introduction), with a rule that a code
change invalidating a document is incomplete until the document is updated in the same commit. D7
records every significant decision (\emph{DEC-nnn}) and every defect found (\emph{F-nn}) with the
reason behind it, and D8 traces each requirement to the code and the tests that verify it. These
records are the source of the figures reported in Chapter~4.

\end{document}
